\section{Filtern}

'Filtern' bedeutet, einen \textbf{optimalen Schätzer} zu finden


\subsection{Optimale Schätzung}

$X_1$ und $X_2$ sind  unabhängige Messungen von $\mu$ mit Messfehler $\var(X_i) = \sigma_i^2 $ 


\subsubsection{Ansatz}

Linearkombination 
$$ \boxed{ X = t \, X_1 + (1-t) \, X_2} $$
so wählen, dass $\var(X)$ minimal wird \\
\textrightarrow\ $X$ entspricht dem 'gewichteten Mittelwert' der Messungen


\subsubsection{Parameter}

\begin{minipage}{0.48\columnwidth}
    $$ \boxed{ t = \frac{\sigma_2^2}{\sigma_1^2 + \sigma_2^2} }$$
\end{minipage}
\hfill
\begin{minipage}{0.48\columnwidth}
    $$ \boxed{ 1-t  = \frac{\sigma_1^2}{\sigma_1^2 + \sigma_2^2} }$$
\end{minipage}


\subsubsection{Lösung}

\begin{minipage}{0.6\columnwidth}
    $$ \boxed{ X = \underbrace{ \frac{\sigma_2^2}{\sigma_1^2 + \sigma_2^2} }_{\substack{t}}  \, X_1 +  
    \underbrace{ \frac{\sigma_1^2}{\sigma_1^2 + \sigma_2^2} }_{\substack{(1-t)}} \, X_2 } $$
\end{minipage}
\hfill
\begin{minipage}{0.38\columnwidth}
    $$ \boxed{ \var(X) = \frac{\sigma_1^2 \, \sigma_2^2}{\sigma_1^2 + \sigma_2^2} } $$
\end{minipage}

